\section{Methodology}
\label{sec:methodology}

The method combines independent onboard recordings to reconstruct a collision
retrospectively. Each recording vehicle, or recorder, contributes a GNSS/INS
pose stream with declared accuracy, control and indicator logs, radar
detections, and collision reports, all on its own clock. Radar supplies
range, azimuth, elevation, and radial velocity without identity, while an
impact report supplies contact time and impulse magnitude without naming
the partner. Each recorder also knows its footprint and sensor mounts; the
investigation supplies the road map and speed limit. Radar observations
from the recorders provide the evidence for reconstructing road users
without accessible recordings. The four stages of
Fig.~\ref{fig:pipeline} read only the inputs shown there.

\begin{figure}[!t]
  \centering
  \begin{tikzpicture}[x=1cm,y=1cm,>=Latex,
  every node/.style={font=\scriptsize},
  stage/.style={draw=black!70, rounded corners=1.5pt, fill=black!5,
    text width=3.55cm, minimum height=0.66cm, align=left, inner sep=2.5pt},
  reads/.style={text width=4.15cm, align=left, inner sep=1pt, text=black!80},
  flow/.style={->, line width=.6pt, draw=black!75}]
  % Inputs available to the method; the simulator's ground truth is not among them.
  \node[draw=black!70, rounded corners=1.5pt, fill=blue!6, text width=8.05cm,
        align=left, inner sep=2.5pt, anchor=north west] (in) at (0.1,0)
    {\textbf{Recordings} of each recorder, on its own clock and frame:
     GNSS/INS pose, radar detections, collision reports, controls, lights,
     footprint. \\ \textbf{Incident context}: road map, speed limit.};
  \node[stage, anchor=north west] (A) at (0.1,-1.30)
    {\textbf{A.} Local reconstruction\\ Kalman/RTS tracks, local events};
  \node[stage, anchor=north west] (B) at (0.1,-2.30)
    {\textbf{B.} Temporal alignment\\ impacts or observed motion};
  \node[stage, anchor=north west] (C) at (0.1,-3.30)
    {\textbf{C.} Identity association\\ recorder links, fused road users};
  \node[stage, anchor=north west] (D) at (0.1,-4.30)
    {\textbf{D.} Shared chronology\\ local events and global identities};
  \node[reads, anchor=west] at (3.95,-1.64)
    {its radar, pose, pedals, lights, collision reports, road map};
  \node[reads, anchor=west] at (3.95,-2.64)
    {collision reports, poses, footprints of all recorders, local tracks};
  \node[reads, anchor=west] at (3.95,-3.64)
    {aligned local tracks, poses with declared accuracy};
  \node[reads, anchor=west] at (3.95,-4.64)
    {local events, matched collisions, identities};
  \draw[flow] (in.south -| A.north) -- (A.north);
  \draw[flow] (A) -- (B);
  \draw[flow] (B) -- (C);
  \draw[flow] (C) -- (D);
\end{tikzpicture}

  \caption{Reconstruction stages and inputs. Local events are extracted
  independently before temporal and identity integration.}
  \label{fig:pipeline}
\end{figure}


\subsection{Local Reconstruction}
\label{subsec:local-reconstruction}

Local processing retains each recorder's clock and uses a frame anchored
at its initial pose. Within this frame, sensor mounting geometry and
recorded ego-motion transform radar positions and compensate Doppler,
including the radar mount's lever arm. Tracks follow the constant-velocity
model of~\eqref{eq:state-space} and are updated with the median position of
associated returns and their Doppler measurements. Before a detection can
update a track, it must pass both spatial and radial-velocity gates. The
position innovation uses~\eqref{eq:association-distance}, with covariance
\begin{equation}
  \begin{aligned}
    \mathbf{S} &= \mathbf{P}_{pp}+2\mathbf{B}+\mathbf{R}(\phi)\operatorname{diag}
      (\sigma_\rho^2,\rho^2\sigma_\phi^2)\mathbf{R}(\phi)^{\mathsf T},\\
    \mathbf{B} &= \mathbf{R}(\hat\theta)\operatorname{diag}
      (L^2/12,W^2/12)\mathbf{R}(\hat\theta)^{\mathsf T},
  \end{aligned}
  \label{eq:gate}
\end{equation}
where $\mathbf{P}_{pp}$ is the predicted position covariance and
$\mathbf{R}(\alpha)$ is a planar rotation. The radar term depends on range
$\rho$, bearing $\phi$, and their measurement standard deviations, while
$\mathbf{B}$ models a uniform point over a nominal $L\times W$ vehicle at
heading $\hat\theta$, with isotropic spread when heading is unavailable.
The factor two accounts for the return and tracked point lying at different
positions on that body. Consequently, the gate widens laterally with range,
while Doppler constrains only line-of-sight velocity.

When several tracks pass both gates, their spatial and Doppler likelihoods
decide the assignment. A detection without a clear likelihood margin
updates no track and starts none, avoiding a forced association. New
tracks require sustained support and are suppressed within a
confirmed body's region. Track continuation ends when prediction
uncertainty prevents separation of vehicle-width bodies, subject to a gap
cap. To accommodate abrupt manoeuvres and impacts, the affected tracks'
motion uncertainty is increased. RTS smoothing~\cite{rauch1965smoothing}
is then applied within segments separated by gaps or detected impacts,
preserving velocity discontinuities.

\subsection{Temporal Alignment}
\label{subsec:temporal-alignment}

Temporal alignment brings these independent local reconstructions onto
a common clock. Collisions provide the main anchors because a contact
between two recorders appears in both reports, with consecutive sensor
callbacks grouped into local contact episodes. Reports from different
recorders are matched on impulse magnitude: the two bodies receive equal
and opposite impulses, so the magnitudes must agree within their
measurement uncertainty. Impulse agreement alone, however, cannot
distinguish unrelated contacts of similar strength. A candidate match
must therefore also satisfy geometric consistency: placed by their own
pose streams, the two footprints must touch within the declared
positioning uncertainty. Where available, pre-contact radar motion must
not contradict the candidate partner's recorded speed, and every
candidate must agree with the clock offsets already established.
If these checks still permit more than one pairing, a match is accepted
only when every feasible one-to-one assignment with the greatest number
of matches contains it. Each accepted match constrains relative clock
offsets, while chains of matched collisions extend the alignment to
recorders that never collided with each other. The strongest matched
collision defines the global time origin, giving
$t_{\mathrm{g}}=t+o_R$ for recorder $R$, where $o_R$ is its
estimated constant offset. The constant-offset model assumes recorder
clocks disciplined by GNSS time, so that drift within one incident can
be neglected.

For a recorder without a shared collision, radar observations of
recording vehicles may still provide alignment. Candidate offsets are
scored by comparing the vehicle surfaces seen by radar with the footprint
and trajectory of the corresponding recorder, within the measurement
uncertainty. The best fit is accepted only when it is consistent and
distinguishable from competing offset intervals; otherwise the recorder
remains unaligned. When no recorder shares a collision at all, one
recorder's clock is taken as the reference.

\subsection{Identity Association}
\label{subsec:identity-association}

With the recorders on a common clock, their tracks can be compared
against the known trajectories of the recorders themselves. These are
the first association candidates; each radar track keeps its local
identifier until it is associated with a road user. A track observed by
recorder $X$ is associated with recorder $Y$ when its reconstructed
position and line-of-sight velocity agree with $Y$'s recorded trajectory
and vehicle footprint. This comparison accounts for radar noise, the
declared positioning accuracy of both recorders, and the displacement
that a heading error causes at range. Linking requires agreement that
lasts long enough and is precise enough to distinguish nearby vehicles,
with a clear likelihood margin when two matches compete. Tracks
approaching a matched collision provide further association evidence.
When positioning uncertainty is declared, observations of recorder
footprints refine relative translation offsets before association and
fusion are repeated; unconstrained components retain their positioning
priors.

Road users without a recorder have no trajectory of their own against
which observations can be checked, so reconstruction instead associates
the tracks of different observers with each other. Direct comparison of
their tracked points would be misleading because each observer sees a
different face of the body. The method therefore tests whether one
vehicle body explains both sets of radar returns. For an estimated
heading, returns are assigned to faces visible from their radar origins;
constraints normal to and along each face then enter a weighted
least-squares fit of the centre, length and width. Measurement and pose
uncertainty determine the constraint weights, while nominal dimensions
provide size priors. Face assignment and fitting are performed iteratively.
The $\chi^2$ terms
below are the resulting weighted residual sums, including the dimension
priors. At each shared observation, the compatibility cost is
\begin{equation}
  D_{\mathrm{body}}=\max\!\left(0,
    \chi^2_{\mathrm{joint}}-\chi^2_a-\chi^2_b\right),
  \label{eq:body-compatibility}
\end{equation}
where the terms are the normalized fitting costs for one joint body and
two separate bodies. This cost is tested against a $\chi^2_2$ gate,
together with velocity compatibility; after either observer's impact,
geometry alone is used. Agreement must persist, and a link is refused
when a rival remains plausible or a \emph{cannot-link} constraint would
be violated: tracks that one observer simultaneously distinguishes as
separate bodies cannot be fused.

Across short gaps, track continuation is evaluated against rival tracks
and a new-road-user hypothesis. Retrospective resolution requires a later
group supported by sustained cross-observer evidence: its first measured
position and velocity are propagated backward with the motion model and
compared with earlier fragments' last measured states, including process
and relative-pose uncertainty. A candidate must pass a $\chi^2_4$
state-compatibility gate and reach model posterior support of at least
0.95 against both competing successors and competing predecessors, with
a new-road-user alternative in each comparison. Merging also requires
that the cannot-link constraints hold and that the earlier fragment ends
without a confusable neighbour. Successful resolution updates identity
attribution while preserving the original observations and event times,
and records when the resolving evidence became available.

Tracks not linked to a recorder form anonymous global road users
$U_1,U_2,\ldots$. Their positions come from the body fitted to the
observed surfaces, and the velocity estimates of different observers
are combined by covariance intersection~\cite{julier1997intersection},
minimizing the trace of the fused covariance. This handles unknown
cross-correlations and also provides the position fallback when no joint
body fit is available. Fused trajectories are smoothed without crossing
gaps or impacts.

\subsection{Semantic Events and Shared Chronology}
\label{subsec:semantic-events}

As part of local reconstruction, semantic events are extracted
independently from each recorder's tracks, motion, controls, lights, and
road context, ready for integration on the recovered clock in the final
stage. Pedal and light logs provide braking and indicator events
directly, whereas interactions are judged against the map-matched path
of the recorder, following lane curvature and its recorded branch at
junctions. An \texttt{EGO\_PATH\_ENTRY} marks
the instant at which the estimated body of another road user enters
this corridor. Since occupancy is computed from the reconstructed
footprint with its position and heading uncertainty, a transition
needs persistent evidence beyond that uncertainty rather than a
single observation. Each entry is then attributed to the target's
lateral motion, the recorder's own manoeuvre, or both, and this
attribution separates a cut-in from a lane change of the recorder:
a \texttt{CUT\_IN} needs a sustained lateral approach from another
lane and is confirmed only by a path entry to which the target's
motion contributes.

Alongside path entries and cut-ins, which describe where other road users
are, a \texttt{CRITICAL\_TTC} event marks when such a situation becomes
critical through either of two triggering conditions. The first is a predicted
overlap between the recorder's safety envelope and another vehicle's
estimated body that cannot be avoided with the assumed deceleration
after a response delay~\cite{green2000stop}. Avoidance is evaluated by
braking the recorder, or the target when recorder braking cannot remove
the overlap. Predictions account for road curvature and estimated target
braking. The second is an unsafe forward gap below a speed-dependent
following-distance reference derived from the time-gap table of UN
Regulation No.~157~\cite{unece2021r157}. This condition applies only to a
leader in the recorder's path or a confirmed cut-in, never to
traffic in adjacent lanes, and unlike a predicted overlap, a short
following gap does not by itself give a
time-to-collision~\cite{vogel2003headway}. Both conditions require
sufficiently mature tracks with bounded position and velocity uncertainty.
The experiments use a 1.0~s response delay, 6.0~m/s$^2$ available
deceleration and a 6.0~s prediction horizon.

Fused trajectories describe the estimated physical motion, while the
shared chronology integrates events already extracted independently by
each recorder. Local semantic events, matched collisions, and associated
identities form a graph on the global clock; each node keeps its recorder,
local event, local timestamp, and provenance. A matched pair of collision
reports becomes one shared node with both participants, while an
unmatched report stays a node of its own recorder. Brief recontacts of
the same pair are grouped with the preceding collision when the two
footprints barely separated in between. Temporal edges then order events
with different reconstructed timestamps, leaving events that share a
timestamp unordered.
